% A number may be assembled out of conditionals, and every part of it is read
% *expanding* — `tex.web` §445 for the signs, §444 for the integer digits, and
% §452 for the decimal point and the digits after it.
%
% The last of those we did not do. The fraction was read with `get_token`, so a
% `\fi` in the middle of it ended the number and everything after the `\fi` was
% left in the input. trip.tex writes
%
%     \scriptspace\if00-0.\fi01pt        % this boils down to -0.01pt
%
% and we read `-0.` and stopped: `\scriptspace` came out `-0.0pt`, which prints
% as `0.0pt`, and the `01` was *typeset* — which started a paragraph, fired
% `\everypar`, and sent the run somewhere TeX never goes. Two of trip.tex's
% spurious pages came from that one missing expansion.
%
% The forms here are the ones a number can be spelled with: a sign, an integer
% digit, a decimal point and a fraction digit, each arriving from a conditional
% rather than from the file.
\catcode`\{=1 \catcode`\}=2 \catcode`\#=6

% The fraction digits come from after a `\fi`.
\dimen0=\if00-0.\fi01pt
\message{<frac-after-fi=\the\dimen0>}

% The decimal point itself arrives from a conditional.
\dimen1=1\if00.\fi5pt
\message{<point-from-if=\the\dimen1>}

% The sign does.
\dimen2=\if00-\fi2pt
\message{<sign-from-if=\the\dimen2>}

% The integer digits do (§444, which we already did).
\dimen3=\if001\fi2pt
\message{<int-from-if=\the\dimen3>}

% The *unit* does — `scan_keyword` expands, so `p` and `t` may arrive apart.
\dimen4=3p\if00t\fi
\message{<unit-from-if=\the\dimen4>}

% A comma is a decimal point too (§452's `continental_point_token`), and it
% works the same way.
\dimen5=\if00-0,\fi25pt
\message{<comma-fraction=\the\dimen5>}

% All of it at once, which is the shape trip.tex actually uses.
\dimen6=\if00-\fi\if000\if00.\fi\fi0\if001\fi p\if00t\fi
\message{<all-at-once=\the\dimen6>}

% Integers scan the same way, including the octal and hex prefixes.
\count0=\ifnum10=10 12\fi
\message{<int=\the\count0>}
\count1='\ifnum10=10 12\fi
\message{<octal=\the\count1>}
\count2="\ifnum10=10 12\fi
\message{<hex=\the\count2>}

\end
